% \documentclass[polish, 8pt]{article}
\documentclass[polish, 10pt]{extarticle}
\usepackage[top=2cm, bottom=2cm, left=3cm, right=3cm]{geometry}
\usepackage{graphicx} % Required for inserting images

\usepackage[polish]{babel}
\usepackage{polski}
\usepackage{placeins}
\usepackage{dirtytalk}
% Set page size and margins
% Replace `letterpaper' with `a4paper' for UK/EU standard size
\usepackage{csquotes}
% Useful packages
\usepackage{float}
\usepackage{amsmath}
\usepackage{graphicx}
\usepackage{hyperref}
\usepackage{graphicx}
\usepackage{listings}
\usepackage{xcolor}
\usepackage{amssymb}
\usepackage{enumitem}
\usepackage{float}
\definecolor{codegreen}{rgb}{0,0.6,0}
\definecolor{codegray}{rgb}{0.5,0.5,0.5}
\definecolor{codepurple}{rgb}{0.58,0,0.82}
\definecolor{backcolour}{rgb}{0.95,0.95,0.92}

\lstdefinestyle{mystyle}{
    backgroundcolor=\color{backcolour},   
    commentstyle=\color{codegreen},
    keywordstyle=\color{magenta},
    numberstyle=\tiny\color{codegray},
    stringstyle=\color{codepurple},
    basicstyle=\ttfamily\footnotesize,
    breakatwhitespace=false,         
    breaklines=true,                 
    captionpos=b,                    
    keepspaces=true,                 
    numbers=left,                    
    numbersep=5pt,                  
    showspaces=false,                
    showstringspaces=false,
    showtabs=false,                  
    tabsize=2
}
\lstset{style=mystyle}
\hypersetup{
    colorlinks=true,
    linktoc=all     %set to all if you want both sections and subsections linked
}
\usepackage{mdframed}
\usepackage{titling}
\setlength{\droptitle}{-2cm} 

\begin{document}

\title{Kolokwium I -- 2. Termin}
\author{}
\date{}

\maketitle
\vspace{-2cm}

\section*{Zadanie 1 (20 pkt)}
W miejscu obsługi podróżnych (MOP) zlokalizowanym przy autostradzie znajduje się stacja benzynowa. Spośród zjeżdżających z autostrady do MOP pojazdów średnio $\frac{1}{4}$ samochodów osobowych, $\frac{2}{8}$ autobusów i $\frac{1}{2}$ samochodów ciężarowych zatrzymuje się na tej stacji, by zatankować paliwo. Wiadomo, że na $100$ zjeżdżających do MOP pojazdów $70$ stanowią samochody osobowe, $10$ autobusy i $20$ samochody ciężarowe. Oblicz prawdopodobieństwo, że pojazd, który został zatankowany na stacji nie jest samochodem osobowym.

\section*{Zadanie 2 (20 pkt)}
Załóżmy, że $P(A) > 0$ i $P(B) > 0$. Udowodnij lub podaj kontrprzykład:

\begin{enumerate}
    \item[(a)] Jeśli $P(A \mid B) > P(A)$, to $P(B \mid A) > P(B)$.
    \item[(b)] Jeśli $P(A \mid B) = P(B \mid A)$, to $P(A) = P(B)$.
    \item[(c)] $P(A \setminus B) = P(A) - P(B)$.
\end{enumerate}

\section*{Zadanie 3 (10 pkt)}
Podaj wypowiedź i dowód twierdzenia o prawdopodobieństwie całkowitym.

\section*{Zadanie 4 (30 pkt)}
Doświadczenie losowe polega na tym, że rzucamy trzema
kostkami symetrycznymi. Odrzucamy skrajne wyniki i 
jako wynik doświadczenia bierzemy pozostały. Niech zmienna
losowa $M$ opisuje wynik tego doświadczenia.
Wyznacz rozkład zmiennej $M$. Wyznacz dystrybuantę tego rozkładu.
Oblicz wartość oczekiwaną, medianę, kwartyl dolny i górny tego rozkładu.

\section*{Zadanie 5 (20 pkt)}
Zmienna losowa $X$ ma rozkład wykładniczy z parametrem $\lambda = 2$. Wyznacz rozkład zmiennej losowej $Y$ zdefiniowanej jako:
\[
Y = e^{-2X}
\]
Jaki rozkład ma zmienna $Y$?
\end{document}

